\subsection*{Procedimiento y operaciones del simulacro 8}
En el actuador se aplican $F=pA$ y $v=Q/A$ con el área útil de cada cámara; las potencias salen del balance hidráulico. La batería se trata como reserva energética utilizable y el cable como una resistencia en serie. El alcance sigue la integración de Breguet para reactor con $L/D$ y consumo constantes.

\begin{enumerate}

\item Presión del sistema.

\[\fitmath{p=21000000\approx 2.1\times10^{7}\,\mathrm{Pa}}\]

\item Área del pistón.

\[\fitmath{\mathrm{Ap}=\frac{\pi}{4}\cdot {0.05}^{2}\approx 0.001963495\,\mathrm{m^{2}}}\]

\item Área del vástago.

\[\fitmath{\mathrm{Ar}=\frac{\pi}{4}\cdot {0.025}^{2}\approx 0.0004908739\,\mathrm{m^{2}}}\]

\item Fuerza máxima en extensión.

\[\fitmath{\mathrm{Fext}=p\cdot \mathrm{Ap}}\]

\[\fitmath{\mathrm{Fext}=(21000000)\cdot (0.001963495408)\approx 41233.4\,\mathrm{N}}\]

\item Fuerza máxima en retracción, con área anular.

\[\fitmath{\mathrm{Fret}=p\cdot \left(\mathrm{Ap}-\mathrm{Ar}\right)}\]

\[\fitmath{\mathrm{Fret}=(21000000)\cdot \left((0.001963495408)-(0.0004908738521)\right)\approx 30925.05\,\mathrm{N}}\]

\item Caudal.

\[\fitmath{Q=0.001\approx 0.001\,\mathrm{m^{3}/s}}\]

\item Velocidad de extensión.

\[\fitmath{\mathrm{vext}=\frac{Q}{\mathrm{Ap}}}\]

\[\fitmath{\mathrm{vext}=\frac{(0.001)}{(0.001963495408)}\approx 0.5092958\,\mathrm{m/s}}\]

\item Velocidad de retracción para el mismo caudal.

\[\fitmath{\mathrm{vret}=\frac{Q}{\mathrm{Ap}-\mathrm{Ar}}}\]

\[\fitmath{\mathrm{vret}=\frac{(0.001)}{(0.001963495408)-(0.0004908738521)}\approx 0.6790611\,\mathrm{m/s}}\]

\item Potencia hidráulica entregada al actuador.

\[\fitmath{\mathrm{Ph}=p\cdot Q}\]

\[\fitmath{\mathrm{Ph}=(21000000)\cdot (0.001)\approx 21000\,\mathrm{W}}\]

\item Potencia mecánica útil con carga de 30 kN.

\[\fitmath{\mathrm{Pu}=30000\cdot \mathrm{vext}}\]

\[\fitmath{\mathrm{Pu}=30000\cdot (0.5092958179)\approx 15278.87\,\mathrm{W}}\]

\item Potencia disipada como calor en válvulas y fugas.

\[\fitmath{\mathrm{loss}=\mathrm{Ph}-\mathrm{Pu}}\]

\[\fitmath{\mathrm{loss}=(21000)-(15278.87454)\approx 5721.125\,\mathrm{W}}\]

\item Potencia de entrada a la bomba con rendimiento 0.85.

\[\fitmath{\mathrm{Pin}=\frac{\mathrm{Ph}}{0.85}}\]

\[\fitmath{\mathrm{Pin}=\frac{(21000)}{0.85}\approx 24705.88\,\mathrm{W}}\]

\item Calentamiento si 18 kg de fluido absorben el calor sin enfriar.

\[\fitmath{\mathrm{dTdt}=\frac{\mathrm{loss}}{18\cdot 1900}}\]

\[\fitmath{\mathrm{dTdt}=\frac{(5721.125463)}{18\cdot 1900}\approx 0.1672844\,\mathrm{K/s}}\]

\item Autonomía con 80 % de la capacidad utilizable.

\[\fitmath{t=\frac{0.8\cdot 40}{15}\approx 2.133333\,\mathrm{h}}\]

\item Resistencia del cableado de cobre de ida y vuelta.

\[\fitmath{R=\frac{1.72\times10^{-8}\cdot 12}{6\times10^{-6}}\approx 0.0344\,\mathrm{ohm}}\]

\item Caída de tensión en el cable.

\[\fitmath{\mathrm{dV}=15\cdot R}\]

\[\fitmath{\mathrm{dV}=15\cdot (0.0344)\approx 0.516\,\mathrm{V}}\]

\item Pérdida por efecto Joule en el cable.

\[\fitmath{\mathrm{Ploss}={15}^{2}\cdot R}\]

\[\fitmath{\mathrm{Ploss}={15}^{2}\cdot (0.0344)\approx 7.74\,\mathrm{W}}\]

\item Tensión en la carga.

\[\fitmath{\mathrm{Vload}=24-\mathrm{dV}}\]

\[\fitmath{\mathrm{Vload}=24-(0.516)\approx 23.484\,\mathrm{V}}\]

\item Caída relativa en por ciento.

\[\fitmath{\mathrm{frac}=\frac{\mathrm{dV}}{24}\cdot 100}\]

\[\fitmath{\mathrm{frac}=\frac{(0.516)}{24}\cdot 100\approx 2.15\,}\]

\item Velocidad de crucero.

\[\fitmath{V=230\approx 230\,\mathrm{m/s}}\]

\item Consumo específico por unidad de peso de empuje.

\[\fitmath{\mathrm{cw}=1.7\times10^{-5}\cdot 9.80665\approx 0.000166713\,\mathrm{1/s}}\]

\item Eficiencia aerodinámica.

\[\fitmath{\mathrm{LD}=16\approx 16\,}\]

\item Logaritmo de la razón de pesos inicial y final.

\[\fitmath{\mathrm{lnW}=\ln\left(1.25\right)\approx 0.2231436\,}\]

\item Alcance de Breguet.

\[\fitmath{\mathrm{Rbr}=\frac{V}{\mathrm{cw}}\cdot \mathrm{LD}\cdot \mathrm{lnW}}\]

\[\fitmath{\mathrm{Rbr}=\frac{(230)}{(0.00016671305)}\cdot (16)\cdot (0.2231435513)\approx 4925639\,\mathrm{m}}\]

\item Autonomía de Breguet.

\[\fitmath{\mathrm{Eh}=\frac{\frac{\mathrm{LD}\cdot \mathrm{lnW}}{\mathrm{cw}}}{3600}}\]

\[\fitmath{\mathrm{Eh}=\frac{\frac{(16)\cdot (0.2231435513)}{(0.00016671305)}}{3600}\approx 5.948839\,\mathrm{h}}\]

\item Fracción de combustible respecto al peso inicial.

\[\fitmath{\mathrm{ff}=1-\frac{1}{1.25}\approx 0.2\,}\]

\end{enumerate}
